\chapter{Modèles profonds en coupe transversale}\label{ch:ml:cross-sectional-deep-models}

Un modèle à facteurs statistiques extrait trois facteurs d'un panel de rendements d'actions et attribue à chaque action une sensibilité fixe à
chacun. Une action cotée la veille n'a pas de sensibilité~; une action dont les caractéristiques ont changé conserve la sensibilité de l'action
qu'elle était. Sur un panel simulé dont les vrais bêtas sont des fonctions de dix caractéristiques, dont la moitié de leur variation est non linéaire,
les sensibilités de l'analyse en composantes principales (ACP) expliquent 22.5~\% des rendements des mois de test, et leurs rendements espérés
sont corrélés à 0.03 avec les vrais. Des bêtas écrits comme fonctions linéaires des caractéristiques (IPCA)
expliquent 27.6~\% et sont corrélés à 0.71~; des bêtas exprimés comme un petit \omterm{def:ml:neural-networks-for-noisy-tabular-data:mlp}{réseau de neurones} appliqué à ces caractéristiques (un \omterm{def:ml:cross-sectional-deep-models:ae}{autoencodeur conditionnel})
expliquent 30.4~\% et sont corrélés à 0.95, contre un plafond de 31.3~\%. Ce chapitre construit des modèles qui apprennent à partir d'un panel
en tant que panel~: poolés entre les actifs, avec des entrées qui décrivent l'actif plutôt que de le désigner.

\section{Panels, poolage et invariance par permutation}

\begin{definition}[Modèle de panel poolé, réseau invariant par permutation]\label{def:ml:cross-sectional-deep-models:pooled}
Un \emph{modèle de panel poolé}\index{modèle de panel poolé} (pooled panel model) ajuste une seule fonction $f_\theta(x_{it})$ à chaque actif $i$ et à chaque date $t$
d'un panel, de sorte que les actifs partagent leurs paramètres à travers leurs caractéristiques plutôt que d'en avoir chacun en propre. Un
\emph{réseau invariant par permutation}\index{réseau invariant par permutation} (permutation-invariant network) associe à un ensemble de vecteurs une sortie
qui ne dépend pas de leur ordre, typiquement $\rho(\sum_i\phi(x_i))$ pour des réseaux $\phi$ et $\rho$ (Zaheer et al., 2017).
\end{definition}

Chaque modèle des chapitres~4 à~7 était poolé~: une seule régression ridge ou un seul réseau pour les 500 actions. Le poolage est ce qui fait
d'un panel de 500 actions et de 240 mois un grand jeu de données~; son prix est l'hypothèse selon laquelle les actions
ayant les mêmes caractéristiques se comportent de manière semblable, ce que dit précisément un modèle factoriel à bêtas pilotés par les caractéristiques.
L'invariance par permutation apparaît dès qu'un modèle a besoin de la coupe transversale entière en entrée~: la moyenne de n'importe quelle caractéristique sur
les actions, ou le rendement d'un portefeuille pondéré par une caractéristique, ne dépend pas de la manière dont les actions sont ordonnées, et c'est
la façon dont les modèles présentés plus bas transforment une coupe transversale de taille quelconque en un nombre fixe de facteurs.

\section{Plongements pour les identifiants et les catégories}

\begin{definition}[Plongement d'entité]\label{def:ml:cross-sectional-deep-models:embedding}
Un \emph{plongement d'entité}\index{plongement d'entité} (entity embedding) représente chaque valeur $j$ d'une variable catégorielle (un secteur, une
plateforme, un identifiant d'action) par un vecteur $e_j\in\mathbb R^m$ appris avec le reste du réseau, à la place d'une colonne one-hot
par catégorie (Guo et Berkhahn, 2016).
\end{definition}

Trente secteurs sont introduits artificiellement dans le panel du chapitre, chacun avec une prime mensuelle tirée avec un écart-type
de 0.4~\%. Un réseau poolé sur les seules dix caractéristiques explique 0.004~\% des rendements centrés des mois de test~; avec un
plongement unidimensionnel du secteur, il explique 0.15~\%, et la corrélation de ses prévisions
avec les vrais rendements espérés passe de 0.52 à 0.75, et le plongement appris est corrélé à 0.86 avec les
primes introduites (les rendements moyens bruts des secteurs sur les mois d'entraînement sont corrélés à 0.88). Un détail pratique
a décidé du résultat~: les plongements initialisés avec l'écart-type par défaut de PyTorch, égal à un, ne se sont jamais assez éloignés
de leur point de départ aléatoire pour que les primes apparaissent~; initialisés près de zéro (écart-type 0.01), ils les ont appris.

\section{Modèles factoriels appris~: analyse en composantes principales instrumentée et autoencodeurs conditionnels}

Le panel du chapitre est \texttt{firm.mlsynth.factor\_panel}~: 300 actions, 240 mois, dix caractéristiques classées, trois
facteurs dont les primes mensuelles sont de 0.6~\%, 0.4~\% et 0.3~\%, et des bêtas qui sont des fonctions des caractéristiques, dont la moitié de leur
variance transversale est linéaire et l'autre moitié non linéaire (carrés, produits, valeurs absolues). Le rendement espéré vrai est
$\mu_{it} = \beta(z_{it})^\top\lambda$.

\begin{definition}[Analyse en composantes principales instrumentée]\label{def:ml:cross-sectional-deep-models:ipca}
L'\emph{analyse en composantes principales instrumentée}\index{analyse en composantes principales instrumentée} (instrumented principal component analysis, IPCA) modélise
$r_{i,t+1} = z_{it}^\top\Gamma f_{t+1} + e_{i,t+1}$ avec $K$ facteurs latents $f$ et des bêtas $z^\top\Gamma$ linéaires en les
caractéristiques (une constante incluse)~; $\Gamma$ ($L\times K$) et les facteurs sont estimés par moindres carrés alternés,
les facteurs mois par mois comme régressions en coupe transversale des rendements sur $Z_t\Gamma$ (Kelly, Pruitt et Su, 2019).
\end{definition}

\begin{definition}[Autoencodeur, autoencodeur conditionnel]\label{def:ml:cross-sectional-deep-models:ae}
Un \emph{autoencodeur}\index{autoencodeur} (autoencoder) est un réseau entraîné à reproduire son entrée à travers un
goulot d'étranglement de dimension réduite~: un encodeur envoie $x$ vers un code $h$, un décodeur renvoie $h$ vers $\hat x$, et la perte est l'erreur de reconstruction.
Un \emph{autoencodeur conditionnel}\index{autoencodeur conditionnel} (conditional autoencoder) reconstruit la coupe transversale des rendements $r_{t+1}$ comme
$\beta(z_{it})^\top f_{t+1}$, où le code $f_{t+1}$ ($K$ facteurs) est une application linéaire de portefeuilles
gérés selon les caractéristiques $Z_t^\top r_{t+1}/n$ et les bêtas sont un \omterm{def:ml:neural-networks-for-noisy-tabular-data:mlp}{réseau de neurones} appliqué aux caractéristiques (Gu, Kelly et Xiu, 2021).
\end{definition}

Les portefeuilles gérés forment la synthèse invariante par permutation de la coupe transversale~: leur nombre est le
nombre de caractéristiques, quel que soit le nombre d'actions. L'\omterm{def:ml:cross-sectional-deep-models:ae}{autoencodeur} du chapitre dispose d'un réseau de bêtas à deux couches cachées (32
et 16 unités) et une application linéaire pour les facteurs~; il est entraîné sur les mois 0 à 159 avec \omterm{def:ml:trees-and-boosting:early}{arrêt précoce} (early stopping) sur les mois 160 à 179~; l'IPCA
et l'ACP sont ajustées sur les mois 0 à 179~; les trois sont évaluées sur les mois 180 à 239.

\begin{proposition}[Pourquoi l'IPCA est une régression]\label{prop:ml:xsnet:ipca}
Pour $\Gamma$ fixé, les facteurs de moindres carrés du mois $t$ valent $\hat f_{t+1} = (\Gamma^\top Z_t^\top Z_t\Gamma)^{-1}\Gamma^\top
Z_t^\top r_{t+1}$~; pour des facteurs fixés, $\operatorname{vec}(\Gamma^\top)$ résout les équations normales
$\sum_t(Z_t^\top Z_t\otimes f_{t+1}f_{t+1}^\top)\operatorname{vec}(\Gamma^\top) = \sum_t(Z_t^\top r_{t+1})\otimes f_{t+1}$.
\end{proposition}

\begin{proof}
La première est une régression par moindres carrés ordinaires de $r_{t+1}$ sur les colonnes de $Z_t\Gamma$. Pour la seconde, $z_{it}^\top\Gamma f =
(z_{it}\otimes f)^\top\operatorname{vec}(\Gamma^\top)$, un modèle linéaire en $\operatorname{vec}(\Gamma^\top)$ dont les
équations normales somment $(z\otimes f)(z\otimes f)^\top = zz^\top\otimes ff^\top$ sur les actifs et les mois.
\end{proof}

\section{Évaluation d'un modèle factoriel appris}

Deux R-carrés mesurent un modèle factoriel (Kelly, Pruitt et Su~; Gu, Kelly et Xiu). Le R-carré \emph{total} utilise les
facteurs réalisés du mois~: la part de la variation de la coupe transversale que les bêtas peuvent expliquer une fois les facteurs connus.
Le R-carré \emph{prédictif} utilise le facteur moyen des mois d'entraînement, $\beta^\top\hat\lambda$~: une
prévision de rendement espéré. Les deux sont mesurés par rapport à zéro.

\begin{table}[htbp]
\centering\small
\begin{tabular}{@{}lccccc@{}}
\toprule
 & \multicolumn{3}{c}{total $R^2$} & \multicolumn{2}{c}{3 facteurs~: corr.\ de $\hat\mu$ avec}\\
\cmidrule(lr){2-4}\cmidrule(l){5-6}
modèle & 1 facteur & 3 facteurs & 5 facteurs & le vrai $\mu$ & partie non linéaire de $\mu$\\
\midrule
ACP (bêtas statiques) & 22.0~\% & 22.5~\% & 23.0~\% & 0.03 & 0.00\\
IPCA (bêtas linéaires) & 25.1~\% & 27.6~\% & 28.1~\% & 0.71 & $-0.03$\\
\omterm{def:ml:cross-sectional-deep-models:ae}{autoencodeur conditionnel} & 26.9~\% & 30.4~\% & 30.8~\% & 0.95 & 0.59\\
\midrule
le vrai & \multicolumn{3}{c}{31.3~\%} & 1 & \\
\bottomrule
\end{tabular}
\caption{Trois modèles factoriels sur les soixante mois de test du panel factoriel conditionnel. Le R-carré total utilise les
facteurs réalisés~; les deux dernières colonnes mesurent la corrélation des rendements espérés de chaque modèle $\beta^\top\hat\lambda$ avec le vrai $\mu$ et avec
la partie de $\mu$ qui provient des termes non linéaires des bêtas. \omterm{def:ml:cross-sectional-deep-models:ae}{Autoencodeur}~: graine~1~; les graines 2 et 3 donnent 30.5~\% et 30.6~\%
de R-carré total, et des corrélations de 0.95 et 0.96. Données~: \texttt{ml\_xsnet}.}\label{tab:ml:xsnet:table}
\end{table}

Le \Cref{tab:ml:xsnet:table} fait apparaître trois constats~: les bêtas statiques ne conviennent pas à un panel dont les caractéristiques évoluent~:
les rendements espérés de l'ACP ne portent aucune information en coupe transversale (corrélation 0.03), bien que son R-carré total, tiré par
le facteur de marché, reste honorable. Les bêtas linéaires retrouvent la partie linéaire du vrai et aucune de la partie non linéaire
($-0.03$). L'\omterm{def:ml:cross-sectional-deep-models:ae}{autoencodeur} retrouve les deux, en grande partie (0.95 au total, 0.59 de la partie non linéaire), et son
R-carré total est à moins d'un point de celui du vrai avec trois facteurs. Le R-carré prédictif, 1.01~\% pour l'ACP, 1.11~\% pour l'IPCA et
1.23~\% pour l'\omterm{def:ml:cross-sectional-deep-models:ae}{autoencodeur}, contre 1.09~\% pour le vrai lui-même, classe les modèles dans le même ordre, mais ne doit pas être lu
isolément~: sur soixante mois, il dépend surtout du niveau des primes et de la moyenne que les rendements des facteurs de la
fenêtre de test ont fini par prendre, ce qui explique que le vrai puisse être dépassé sur ce critère.

\begin{omfigure}
\begin{tikzpicture}
\begin{axis}[width=10.5cm,height=5.4cm,xlabel={nombre de facteurs $K$},ylabel={$R^2$ total hors échantillon (\%)},
  tick label style={font=\small},label style={font=\small},xmin=0.5,xmax=5.5,ymin=20,ymax=33,xtick={1,3,5},grid=major,
  grid style={gray!25},legend columns=4,legend style={font=\footnotesize,at={(0.5,-0.3)},anchor=north,draw=none,
  fill=none,/tikz/every even column/.append style={column sep=6pt}},table/col sep=comma]
\addplot[gray,very thick,mark=triangle*] table[x=K,y=pca]{figdata/ml/09-cross-sectional-deep-models/total.csv};
\addplot[omThm,very thick,mark=square*] table[x=K,y=ipca]{figdata/ml/09-cross-sectional-deep-models/total.csv};
\addplot[omDef,very thick,mark=*] table[x=K,y=cae]{figdata/ml/09-cross-sectional-deep-models/total.csv};
\addplot[black!70,thick,dashed] table[x=K,y=truth]{figdata/ml/09-cross-sectional-deep-models/total.csv};
\legend{PCA,IPCA,autoencodeur conditionnel,vrai modèle (3 facteurs)}
\end{axis}
\end{tikzpicture}

\omcaption{R-carré total des mois de test en fonction du nombre de facteurs. Le vrai modèle compte trois facteurs~; les modèles qui en ont davantage
ont plus de marge d'ajustement, et non davantage de vérité à retrouver. Données~: \texttt{ml\_xsnet.pca}, \texttt{ipca}, \texttt{cae}.}\label{fig:ml:xsnet:total}
\end{omfigure}

\begin{method}[Ajuster un modèle factoriel appris]\label{met:ml:xsnet:method}
\begin{enumerate}
\item Classer les caractéristiques par date et ajouter une constante~; conserver le rendement du mois $t+1$ avec les caractéristiques du
mois $t$.
\item Ajuster d'abord l'IPCA~: elle est rapide, n'exige aucune graine et indique la part qu'explique un modèle à bêtas linéaires.
\item Ajuster l'\omterm{def:ml:cross-sectional-deep-models:ae}{autoencodeur conditionnel} avec \omterm{def:ml:trees-and-boosting:early}{arrêt précoce} sur des mois de validation et plusieurs graines~; comparer à l'IPCA sur
le R-carré total et sur la corrélation des rendements espérés avec une prévision de référence.
\item Choisir le nombre de facteurs sur des mois de validation~; rapporter les deux R-carrés et ne jamais se fonder sur le R-carré prédictif seul.
\end{enumerate}
\end{method}

\section{Tutoriel~: des facteurs qui apprennent}\label{tut:ml:cross-sectional-deep-models:tut}

\begin{tutorial}
\textbf{Objectif.} Ajuster l'ACP, l'IPCA et un \omterm{def:ml:cross-sectional-deep-models:ae}{autoencodeur conditionnel} à un, trois et cinq facteurs sur un panel dont les bêtas
dépendent des caractéristiques~; mesurer la part de non-linéarité injectée que chacun retrouve~; ajouter les secteurs par un
plongement. \textbf{État final~:} le \cref{tab:ml:xsnet:table}, la \cref{fig:ml:xsnet:total}.
\begin{tutsteps}
\item \textbf{IPCA} par moindres carrés alternés (la \cref{prop:ml:xsnet:ipca}).
\omcode{code/firm/xsnet/firm_xsnet.py}{66}{89}{Analyse en composantes principales instrumentée~: facteurs par régression en coupe transversale, charges par les équations normales de Kronecker.}\label{lst:ml:xsnet:ipca}
\item \textbf{L'\omterm{def:ml:cross-sectional-deep-models:ae}{autoencodeur conditionnel}}~: bêtas issus des caractéristiques, facteurs issus de portefeuilles gérés.
\omcode{code/firm/xsnet/firm_xsnet.py}{100}{111}{Les deux réseaux de l'autoencodeur conditionnel.}\label{lst:ml:xsnet:ca}
\item \textbf{Lancer} \texttt{ml\_xsnet.pca(K)}, \texttt{ipca(K)}, \texttt{cae(K, seed)}, \texttt{embedding()} et
\texttt{fig\_xsnet.py}.
\end{tutsteps}
\textbf{Pour aller plus loin.} Mettre à zéro la part non linéaire injectée et vérifier que l'IPCA rattrape le
\omterm{def:ml:cross-sectional-deep-models:ae}{autoencodeur}~; donner aussi à la carte des facteurs de l'\omterm{def:ml:cross-sectional-deep-models:ae}{autoencodeur} les portefeuilles gérés des caractéristiques élevées au carré.
\end{tutorial}

\section{Construction~: réseaux en coupe transversale et modèles factoriels}\label{bld:ml:cross-sectional-deep-models:xsnet}

\begin{build}
\bfield{Rôle} Les modèles factoriels appris de la maison, comparables sur un même panel et une même paire de R-carrés.
\bfield{Interface} \texttt{managed(Z, r)}, \texttt{with\_constant(Z)}, \texttt{PCAModel(K)}, \texttt{IPCA(K, iters)}
avec \texttt{fit}, \texttt{factors}, \texttt{betas}, \texttt{score}~; \texttt{CondAE(L, K, hidden, seed)} avec les mêmes
méthodes~; \texttt{scores(pred\_total, pred\_mu, r)}~; \texttt{EmbedMLP}, \texttt{fit\_embed}.
\bfield{Règles} Caractéristiques du mois $t$ avec rendements du mois $t+1$~; portefeuilles gérés calculés sur les caractéristiques et les
rendements d'un même mois~; un seul thread CPU, \omterm{def:ml:trees-and-boosting:early}{arrêt précoce} sur des mois de validation.
\bfield{Tests d'acceptation} \texttt{code/firm/xsnet/tests/}~: l'IPCA retrouve les bêtas d'un modèle à bêtas linéaires injecté (à rotation près)
et ses étapes alternées ne font jamais croître la perte~; l'ACP retrouve les chargements statiques~; l'entrée en portefeuilles gérés
de l'\omterm{def:ml:cross-sectional-deep-models:ae}{autoencodeur} est invariante à la permutation des titres~; un plongement retrouve les effets de catégorie injectés.
\bfield{Pour aller plus loin} Test de l'IPCA avec terme constant (alpha)~; l'\omterm{def:ml:cross-sectional-deep-models:ae}{autoencodeur} avec cartes de facteurs non linéaires~; le critère
d'absence d'arbitrage de Chen, Pelger et Zhu.
\end{build}

\begin{omsources}
\item B.~Kelly, S.~Pruitt et Y.~Su, «~Characteristics are covariances: a unified model of risk and return~», \emph{Journal
of Financial Economics} 134(3), 2019.
\item S.~Gu, B.~Kelly et D.~Xiu, «~Autoencoder asset pricing models~», \emph{Journal of Econometrics} 222(1), 2021.
\item L.~Chen, M.~Pelger et J.~Zhu, «~Deep learning in asset pricing~», \emph{Management Science} 70(2), 2024.
\item C.~Guo et F.~Berkhahn, «~Entity embeddings of categorical variables~», arXiv, 2016.
\item M.~Zaheer et al., «~Deep sets~», \emph{NeurIPS}, 2017.
\end{omsources}

\section{Exercices}

\begin{exercise}[$\star$]\label{exo:ml:cross-sectional-deep-models:1}
Avec dix caractéristiques, une constante et trois facteurs, calculer le nombre de paramètres de $\Gamma$ de l'IPCA, puis celui d'un
modèle ACP statique de 300 titres à trois facteurs.
\end{exercise}

\begin{exercise}[$\star$]\label{exo:ml:cross-sectional-deep-models:2}
Trois titres ont pour caractéristiques 0.5, $-0.5$, 1.0 et pour rendements 2~\%, $-1~\%$, 3~\% sur un mois. Calculer le
rendement du portefeuille géré de la caractéristique, puis indiquer s'il dépend de l'ordre des titres.
\end{exercise}

\begin{exercise}[$\star$]\label{exo:ml:cross-sectional-deep-models:3}
Pour trente secteurs, comparer un codage one-hot injecté dans une couche cachée de 32 unités à un plongement unidimensionnel
alimentant la même couche, et indiquer combien de paramètres chacun utilise pour le secteur.
\end{exercise}

\begin{exercise}[$\star\star$]\label{exo:ml:cross-sectional-deep-models:4}
Expliquer pourquoi le rendement espéré de l'ACP ne porte aucune information en coupe transversale sur ce panel, alors que son R-carré total atteint 22.5~\%.
\end{exercise}

\begin{exercise}[$\star\star$]\label{exo:ml:cross-sectional-deep-models:5}
Expliquer pourquoi un modèle peut battre le vrai sur le R-carré prédictif calculé sur soixante mois, et ce qu'il faut rapporter à la place.
\end{exercise}

\begin{exercise}[$\star\star$]\label{exo:ml:cross-sectional-deep-models:6}
\emph{Trouver la faille.} «~Notre \omterm{def:ml:cross-sectional-deep-models:ae}{autoencodeur conditionnel} à dix facteurs explique 40~\% des rendements dans l'échantillon (in-sample), contre 30~\% pour l'IPCA,
donc il saisit davantage de la coupe transversale.~»
\end{exercise}

\begin{exercise}[$\star\star\star$]\label{exo:ml:cross-sectional-deep-models:7}
\emph{Programmation.} Relancer \texttt{ml\_xsnet.embedding} avec des plongements initialisés à la valeur par défaut de PyTorch (écart type égal à un).
Rapporter la corrélation du plongement avec les primes et expliquer.
\end{exercise}

\begin{exercise}[$\star\star\star$]\label{exo:ml:cross-sectional-deep-models:8}
Montrer que l'IPCA avec $\Gamma$ égal à l'identité et $K = L$ se ramène à la régression des rendements de chaque mois sur ses
caractéristiques (première étape de Fama et MacBeth, livre~4, chapitre~16).
\end{exercise}

\section{Problème~: Facteurs qui apprennent}

\begin{problem}[{Problème du week-end --- bêtas qui varient avec les caractéristiques}]\label{pb:ml:cross-sectional-deep-models:1}
Le panel du chapitre~: 300 actions, 240 mois, dix caractéristiques, trois facteurs, bêtas à moitié non linéaires.

\medskip\noindent\textbf{Partie I --- Le vrai modèle.}
\begin{enumerate}
\item Quels R-carrés total et prédictif le vrai modèle obtient-il sur les mois de test~?
\item Pourquoi le R-carré total du vrai modèle est-il loin de l'unité~?
\item Que sont les primes, et que font-elles au R-carré prédictif~?
\item Que signifie «~à moitié non linéaire~» pour un modèle à bêtas linéaires~?
\end{enumerate}

\medskip\noindent\textbf{Partie II --- Trois modèles.}
\begin{enumerate}[resume]
\item Donner le R-carré total de chaque modèle avec un, trois et cinq facteurs.
\item Dans quelle mesure les rendements espérés de chaque modèle sont-ils corrélés au vrai modèle et à sa partie non linéaire~?
\item Dans quelle mesure les graines aléatoires font-elles varier l'\omterm{def:ml:cross-sectional-deep-models:ae}{autoencodeur}~?
\item Pourquoi l'IPCA ne retrouve-t-elle aucune part de la partie non linéaire~?
\end{enumerate}

\medskip\noindent\textbf{Partie III --- Poolage et plongements.}
\begin{enumerate}[resume]
\item Pourquoi le poolage est-il l'hypothèse appropriée pour ce panel~?
\item Qu'apporte le plongement sectoriel, et à quel point se rapproche-t-il des moyennes sectorielles brutes~?
\item Qu'est-ce qui a déterminé si le plongement a appris quoi que ce soit~?
\item Où l'invariance par permutation apparaît-elle dans l'\omterm{def:ml:cross-sectional-deep-models:ae}{autoencodeur}~?
\end{enumerate}

\medskip\noindent\textbf{Partie IV --- Le verdict.}
\begin{enumerate}[resume]
\item Énoncer le \emph{résultat nommé}~: R-carrés total et prédictif hors échantillon de l'ACP, de l'IPCA et de l'\omterm{def:ml:cross-sectional-deep-models:ae}{autoencodeur} à trois
facteurs, ainsi que la part de la non-linéarité introduite que chacun retrouve.
\item Quel modèle retiendrait-on pour construire un portefeuille neutre aux facteurs, et pourquoi~?
\item Comment choisir le nombre de facteurs~?
\item Qu'est-ce qui change avec de vraies caractéristiques bruitées et lacunaires~?
\item Que donne chaque modèle pour une action cotée la veille~?
\item Quel coût l'\omterm{def:ml:cross-sectional-deep-models:ae}{autoencodeur} impose-t-il que l'IPCA n'impose pas~?
\item Comment tester que l'ajustement supplémentaire de l'\omterm{def:ml:cross-sectional-deep-models:ae}{autoencodeur} ne relève pas du \omterm{def:ml:why-financial-machine-learning-is-different:generalisation}{surapprentissage}~?
\item En une phrase~: qu'apprend un modèle factoriel entraîné~?
\end{enumerate}
\end{problem}

\section{Questions d'entretien}

\begin{interviewq}[$\star$ \iqroles{chercheur}]\label{iq:ml:cross-sectional-deep-models:1}
Quel problème l'IPCA résout-elle que l'ACP sur les rendements ne résout pas~?
\end{interviewq}

\begin{interviewq}[$\star\star$ \iqroles{chercheur, ingénieur ML}]\label{iq:ml:cross-sectional-deep-models:2}
Expliquez les deux réseaux de l'\omterm{def:ml:cross-sectional-deep-models:ae}{autoencodeur conditionnel} et ce qui entre dans chacun.
\end{interviewq}

\begin{interviewq}[$\star\star$ \iqroles{ingénieur ML}]\label{iq:ml:cross-sectional-deep-models:3}
Quand utiliseriez-vous un \omterm{def:ml:cross-sectional-deep-models:embedding}{plongement d'entité} pour un identifiant d'action, et qu'est-ce qui peut mal tourner~?
\end{interviewq}

\begin{interviewq}[$\star\star$ \iqroles{chercheur}]\label{iq:ml:cross-sectional-deep-models:4}
R-carré total ou R-carré prédictif~: lequel retenez-vous pour comparer des modèles factoriels, et pourquoi~?
\end{interviewq}

\begin{interviewq}[$\star\star$ \iqroles{chercheur, ingénieur ML}]\label{iq:ml:cross-sectional-deep-models:5}
Comment alimenteriez-vous un modèle avec l'information sur l'ensemble de la coupe transversale à une date donnée, quel que soit le nombre d'actions~?
\end{interviewq}

\begin{interviewq}[$\star\star\star$ \iqroles{chercheur}]\label{iq:ml:cross-sectional-deep-models:6}
Dérivez les étapes de moindres carrés alternés de l'IPCA.
\end{interviewq}
